\DeclareEditionOverlay{CM-C01-U001}{%
\begin{EditionUnitCard}{Antes de empezar}
Los sistemas numéricos se amplían cuando una operación o una medida exige
objetos que el sistema anterior no contiene.
\RegisteredBookFigure{FIG-C01-001}{fig:c01:number-systems}{Sistemas numéricos
anidados.}{Diagrama de conjuntos anidados: los naturales están contenidos en
los enteros, estos en los racionales y estos en los reales; la zona real
exterior a los racionales contiene irracionales como raíz de dos y pi.}{\NumberSystemsVisual}
\begin{EditionCheckpointBox}{Comprueba}
Sitúa en el diagrama $-2$, $3/5$ y $\sqrt2$. ¿En qué conjuntos está cada uno?
\end{EditionCheckpointBox}
\begin{EditionWarningBox}{No es una jerarquía de importancia}
La inclusión muestra pertenencia. No dice que un tipo de número sea mejor que
otro ni abre aquí un estudio de cardinalidad.
\end{EditionWarningBox}
\end{EditionUnitCard}}

\DeclareEditionOverlay{CM-C01-U002}{%
\begin{EditionUnitCard}{Estrategia de lectura}
Antes de seguir cada igualdad, identifica cuatro momentos: suposición
contraria, representación reducida, consecuencia de paridad y contradicción.
\begin{EditionCheckpointBox}{Recupera la idea}
Sin mirar el texto, di qué afirmación resulta incompatible con que la fracción
esté reducida.
\end{EditionCheckpointBox}
\end{EditionUnitCard}}

\DeclareEditionOverlay{CM-C01-U003}{%
\begin{EditionUnitCard}{Dos ideas distintas}
Densidad habla de puntos entre puntos; completitud, de la existencia de una
frontera real para ciertos conjuntos acotados.
\begin{EditionWarningBox}{Resultado de referencia}
Aquí no se construye $\R$. Estudia qué permite afirmar la completitud, no una
demostración mediante una construcción de los reales.
\end{EditionWarningBox}
\end{EditionUnitCard}}

\DeclareEditionOverlay{CM-C01-U004}{%
\begin{EditionUnitCard}{Ver el valor absoluto como distancia}
La lectura geométrica traduce condiciones algebraicas a regiones de la recta.
\RegisteredBookFigure{FIG-C01-002}{fig:c01:distance-neighbourhood}{Entorno
abierto y distancia al centro.}{Recta real con el entorno abierto de centro a
y radio r; círculos huecos excluyen a menos r y a más r, y un segmento muestra
la distancia absoluta entre x y a.}{\DistanceNeighbourhoodVisual}
\begin{EditionCheckpointBox}{Piénsalo}
Explica por qué los círculos en $a-r$ y $a+r$ son huecos y cómo se traduce ese
detalle en la desigualdad $|x-a|<r$.
\end{EditionCheckpointBox}
\end{EditionUnitCard}}

\DeclareEditionOverlay{CM-C01-U005}{%
\begin{EditionUnitCard}{Tres preguntas diferentes}
\begin{EditionWarningBox}{Conserva la distinción}
¿Es una cota? ¿Pertenece al conjunto? ¿Es la menor cota superior? Responder una
pregunta no responde automáticamente las otras.
\end{EditionWarningBox}
\begin{EditionCheckpointBox}{Antes de seguir}
Compara un intervalo que incluya su extremo derecho con otro que no: ¿qué cambia
para el máximo y qué permanece igual para el supremo?
\end{EditionCheckpointBox}
\end{EditionUnitCard}}

\DeclareEditionOverlay{CM-C01-U006}{%
\begin{EditionUnitCard}{Una función es más que una expresión}
Pregunta siempre qué entradas se admiten, a qué conjunto deben pertenecer las
salidas y qué regla asigna cada valor.
\begin{EditionWarningBox}{Atención al contexto}
El dominio natural de una expresión y el dominio de una variable en un modelo
aplicado pueden ser distintos.
\end{EditionWarningBox}
\end{EditionUnitCard}}

\DeclareEditionOverlay{CM-C01-U007}{%
\begin{EditionUnitCard}{Seguir entradas y salidas}
\RegisteredBookFigure{FIG-C01-003}{fig:c01:function-mapping}{Dominio,
codominio e imagen de una función.}{Dos conjuntos ovalados muestran dominio y
codominio; tres flechas alcanzan solo una parte del codominio y un punto queda
sin imagen inversa.}{\FunctionMappingVisual}
\begin{EditionWarningBox}{Dos usos que no debes confundir}
La preimagen de un conjunto existe para cualquier función. Una función inversa
exige invertibilidad.
\end{EditionWarningBox}
\begin{EditionCheckpointBox}{Antes de seguir}
Localiza un valor del codominio que no esté en la imagen. ¿Qué propiedad falla
por esa razón?
\end{EditionCheckpointBox}
\end{EditionUnitCard}}

\DeclareEditionOverlay{CM-C01-U008}{%
\begin{EditionUnitCard}{Diagnóstico de invertibilidad}
Antes de hallar una inversa, revisa inyectividad y sobreyectividad. Si una falla,
pregunta si una restricción coherente produce una función nueva e invertible.
\begin{EditionWarningBox}{La expresión no decide sola}
Dos funciones con la misma fórmula pueden tener propiedades distintas porque
sus dominios o codominios son diferentes.
\end{EditionWarningBox}
\end{EditionUnitCard}}

\DeclareEditionOverlay{CM-C01-U009}{%
\begin{EditionUnitCard}{Componer con tipos}
Lee de derecha a izquierda y dibuja la cadena de conjuntos por la que pasan los
valores.
\begin{EditionCheckpointBox}{Comprueba}
Antes de calcular, escribe qué valores salen de la primera función y verifica
si la segunda puede recibirlos.
\end{EditionCheckpointBox}
\end{EditionUnitCard}}

\DeclareEditionOverlay{CM-C01-U010}{%
\begin{EditionUnitCard}{Revisión selectiva}
El objetivo es conservar información de dominio mientras simplificas y
reconocer cómo se ensamblan los trozos, no repetir toda el álgebra escolar.
\begin{EditionWarningBox}{Simplificar no cambia el dominio original}
Un punto excluido por el denominador original sigue excluido aunque la forma
simplificada ya no muestre ese factor.
\end{EditionWarningBox}
\end{EditionUnitCard}}

\DeclareEditionOverlay{CM-C01-U011}{%
\begin{EditionUnitCard}{Exponencial y logaritmo}
Estudia estas funciones como una pareja inversa y escribe las restricciones
junto a cada identidad.
\begin{EditionWarningBox}{Notación y dominio}
El logaritmo natural se escribe $\ln$. Un logaritmo real necesita argumento
positivo; una base declarada debe ser positiva y distinta de uno.
\end{EditionWarningBox}
\begin{EditionCheckpointBox}{Recupera las condiciones}
Antes de aplicar una identidad logarítmica, enumera las condiciones de sus
argumentos y de la base.
\end{EditionCheckpointBox}
\end{EditionUnitCard}}

\DeclareEditionOverlay{CM-C01-U012}{%
\begin{EditionUnitCard}{Radianes e inversas}
Relaciona la medida angular con la circunferencia unidad y presta atención a
los puntos donde un cociente trigonométrico no está definido.
\begin{EditionWarningBox}{Las ramas importan}
Una función periódica necesita restringirse antes de tener inversa. Por eso una
identidad inversa puede ser válida solo en una rama.
\end{EditionWarningBox}
\end{EditionUnitCard}}

